\documentclass[11pt]{article}
\usepackage[a4paper,margin=28mm]{geometry}
\usepackage{amsmath,amssymb,amsthm,mathtools,hyperref}
\newcommand{\D}{dd^c}
\newcommand{\PSH}{\operatorname{PSH}}
\newcommand{\MA}{\operatorname{MA}}
\newcommand{\Eenv}{\mathcal E_{\mathrm{env}}}
\newtheorem{theorem}{Theorem}
\newtheorem{lemma}[theorem]{Lemma}
\newtheorem{corollary}[theorem]{Corollary}
\theoremstyle{remark}\newtheorem{remark}[theorem]{Remark}
\title{Residual regularity of singular radial weighted envelopes}
\author{}\date{}
\begin{document}\maketitle
\begin{abstract}
We consider an ordinary global obstacle envelope with a sublogarithmic
pole on projective space. A compactly supported radial perturbation of
the singular weight reduces its envelope to a one-dimensional convex
minorant. The residual against the prescribed singular reference is
globally $C^{1,1}$, vanishes near the pole, and is smooth near the
hyperplane at infinity. The envelope retains its ordinary type and its
actual non-pluripolar measure has mass one, with zero density on every
free interval. Fixed-support bounded obstacle families have a uniform
$C^{1,1}$ residual bound and quantitative $C^{1,\gamma}$ stability for
$0<\gamma<1$, even when the infinity contact moves. The formula extends to the full family of competitors on
a product with an arbitrary positive Hermitian factor, including a
nonclosed background on a non-K\"ahler fourfold. The conclusions concern
this declared singular-obstacle operator and its residual.
\end{abstract}
\noindent\textbf{Checked scoped result.} The declared theorem has been
independently reconstructed and reviewed. The unrestricted original
S06 programme remains open.

\section{The singular reference and the precise operator}
Use $\D=i\partial\bar\partial/(2\pi)$ and ordinary wedge powers.
Let $\kappa=\D\log(1+|z|^2)$ be the standard Fubini--Study form on
$\mathbf P^k$, $k\geq2$, with $\int\kappa^k=1$. On its standard affine
chart put
\[
 r=|z|^2,\qquad x=\log r,\qquad p=\frac r{1+r},\qquad
 T=1-\log p,\qquad 0<\alpha<1,
\]
and define
\[
 \psi=1-T^\alpha,\qquad
 A(x)=\log(1+e^x)+\psi(x).
\]
The canonical value of $\psi$ at $z=0$ is $-\infty$.
At infinity, with $R=1/r$, it is
$1-(1+\log(1+R))^\alpha$, a smooth expression in $R$.
Here $R$ is a smooth projective function near the hyperplane, and the
expression has value zero there. Thus $\psi$ is a genuine global
$\kappa$-psh function, as also follows from the derivative
calculation below. All inequalities involving potentials include their
canonical values. Ordinary order $u\preceq v$ means $u\leq v+C_0$
for a finite constant $C_0$.

Fix $F\in C_c^2(\mathbf R,\mathbf R)$ and pull $F(\log r)$ back to
the punctured chart, extending it by zero at both ends. Define the
ordinary global obstacle envelope
\begin{equation}\label{eq:envelope}
 E_F=\left(\sup\{v\in\PSH(\mathbf P^k,\kappa):
       v\leq\psi+F\text{ everywhere},\ v\preceq\psi\}\right)^*.
\end{equation}
Here the obstacle set is the entire manifold. Since $F$ is bounded,
the order condition is redundant for each competitor; its inclusion
declares the intended ordinary class. This is not an ideal-order
envelope. The weight itself takes the value $-\infty$ at the pole.

\begin{lemma}\label{lem:reference}
$A$ is smooth and strictly convex, $0<A'<1$, and $A''$ is bounded.
Its asymptotic slopes are $0$ and $1$. Moreover,
\[
 \begin{split}
 &A(x)=1-(1-x)^\alpha+o(1)\quad(x\longrightarrow-\infty),\\
 &A(x)-xA'(x)\longrightarrow-\infty\quad(x\longrightarrow-\infty),\\
 &A(x)-x\searrow0,\quad A(x)-xA'(x)\longrightarrow0
       \quad(x\longrightarrow+\infty).
 \end{split}
\]
\end{lemma}
\begin{proof}
Put $a=\alpha T^{\alpha-1}$. Since $p'=p(1-p)$ and $T'=-(1-p)$,
\[
 A'=p+(1-p)a,\qquad
 A''=p(1-p)(1-a)+\alpha(1-\alpha)(1-p)^2T^{\alpha-2}.
\]
Both assertions follow from $T\geq1$ and $0<a\leq\alpha<1$.
The left asymptotics follow from $T=1-x+O(e^x)$; the tangent intercept
has leading term $-(1-\alpha)(-x)^\alpha$.
On the right, $A=x+(1-\alpha)e^{-x}+O(e^{-2x})$ and
$A'=1-(1-\alpha)e^{-x}+O(e^{-2x})$. Finally $A'-1<0$.
\end{proof}

\section{Convex reduction for the full envelope}
Let $g=A+F$ and let $C$ be its largest convex minorant on $\mathbf R$,
defined as the supremum of all affine functions below $g$.
This gives a finite convex function because $A-\|F\|_\infty$ is a
convex minorant and is itself the supremum of its supporting lines.
In particular
\begin{equation}\label{eq:sandwich}
 A-\|F\|_\infty\leq C\leq A+F\leq A+\|F\|_\infty.
\end{equation}

\begin{theorem}\label{th:main}
The exact envelope in \eqref{eq:envelope} is
\[
 E_F(z)=C(\log|z|^2)-\log(1+|z|^2)
\]
off its pole, with canonical extensions at the pole and infinity.
It is ordinarily equivalent to $\psi$. Its residual $E_F-\psi$ has a
canonical extension in $C^{1,1}(\mathbf P^k)$, is identically zero
near the pole, and is smooth near infinity.
If $J=\{C<g\}$, then $C$ is affine on each connected component of $J$,
and, writing $L_2=\sup_{\mathbf R}\max(g'',0)$,
\begin{equation}\label{eq:quantitative}
 0\leq C'(b)-C'(a)\leq L_2(b-a)\qquad(a<b).
\end{equation}
\end{theorem}
\begin{proof}[Proof of the formula and type]
Every supporting slope of $C$ lies in $[0,1]$: a negative slope is
incompatible with $g(x)\to-\infty$ at the left end, and a slope above
one is incompatible with $g(x)=x+o(1)$ at the right end.
The bounds \eqref{eq:sandwich} force the limiting slopes to be $0,1$.

The full competitor family is invariant under $U(k)$ acting in the affine
chart, so its upper-regularized supremum is radial. It is genuine
$\kappa$-psh: local smooth closed potentials give a locally uniformly
upper-bounded psh family, whose upper-regularized supremum is psh.
The competitor $\psi-\|F\|_\infty$ prevents collapse. The obstacle
$\psi+F$ is upper semicontinuous, so regularization preserves the
pointwise upper bound, including the pole.

For completeness, a radial psh total potential has the form $B(\log r)$
with $B$ convex and nondecreasing. Restriction along a complex line gives
convexity in $\log r$: radial subharmonic functions satisfy
$(r\,\partial_r B(\log r))'\geq0$ in distributions.
On the punctured affine space, positivity in a complex direction
orthogonal to $z$ gives $B'\geq0$ for $k\geq2$.
These statements hold for general radial psh functions by invariant
local smoothing on annuli and distributional limits. Conversely, for
smooth convex nondecreasing $B$, the complex Hessian has tangential
eigenvalue $B'/r$ and radial eigenvalue $B''/r$; one-dimensional smooth
convex approximation proves the distributional assertion for finite
convex $B$. The envelope total profile is therefore a convex minorant
of $g$ and is at most $C$.

Conversely $C$ gives a psh total potential on the punctured chart by
this Hessian criterion. Its upper bound at the origin and
\eqref{eq:sandwich} give a canonical removable psh extension with value
$-\infty$. Since $C'\leq1$, $C-x$ is nonincreasing and has a finite
limit at infinity by \eqref{eq:sandwich}. The proposed potential is
locally bounded above there. Adding a local smooth potential of
$\kappa$ and applying removable psh singularities across the
hyperplane proves a genuine extension. Thus it is a global competitor,
and the formula follows. The bounds give ordinary equivalence without
ever subtracting two infinite pole values.
\end{proof}

\begin{proof}[Proof of regularity and the end formulas]
On a compact interval inside $J$ the gap $g-C$ has a positive minimum.
Replacing $C$ on a sufficiently short subinterval by its chord preserves
convex gluing: the chord slope lies between the corresponding endpoint
one-sided derivatives. Continuity makes the increase smaller than the
gap. Unless $C$ is locally affine this gives a larger convex minorant,
a contradiction. Overlapping intervals show that $C$ is affine on
each component of $J$.

At a contact point $y$, domination by the differentiable obstacle gives
$C'_-(y)\geq g'(y)\geq C'_+(y)$; convexity gives the reverse order.
Thus $C$ is differentiable there, and both derivatives equal $g'(y)$.
It is differentiable on free intervals as well. A differentiable convex
function on an open interval has continuous derivative, so $C\in C^1$.
For $a<b$ in distinct free intervals, move $a$ to its right contact
endpoint and $b$ to its left contact endpoint, leaving contact points
fixed. The new points $a',b'$ lie in $[a,b]$ and have the same
derivatives. Any unbounded free component that has no appropriate
endpoint puts both original points in that component, where the
difference is zero. Otherwise
\[
 0\leq C'(b)-C'(a)=g'(b')-g'(a')
       \leq L_2(b'-a')\leq L_2(b-a).
\]
This proves \eqref{eq:quantitative}.

For a tangent $\ell_x(y)=A'(x)y+A(x)-xA'(x)$, Lemma~\ref{lem:reference}
gives $\ell_x\to-\infty$ uniformly on $\operatorname{supp}F$ as
$x\to-\infty$. For all sufficiently negative $x$, this tangent is
below $g$ there, and below $A=g$ outside the support. It is an affine
minorant touching $A$ at $x$, so $C(x)=A(x)$ on a left end interval.
The residual is exactly zero near the pole, and this is its canonical
smooth extension.

Put $L=\inf_{\mathbf R}(g(x)-x)\leq0$. If $L$ is attained at $y$,
the minorant $x+L$ touches $g$ at $y$. The slope bound $C'\leq1$
then forces $C=x+L$ on $[y,\infty)$: indeed
$C(x)\leq C(y)+(x-y)=x+L$ and the reverse inequality holds everywhere.
If $L<0$, it is attained, since $g-x\to+\infty$ on the left and
to zero on the right. The same formula holds if $L=0$ is attained.
The envelope is then
\begin{equation}\label{eq:tail}
 E_F=L-\log(1+R)
\end{equation}
near infinity, a smooth projective function there.

If $L=0$ is not attained, $g(y)-y$ has a positive minimum on
$\operatorname{supp}F$. The tangents $\ell_x(y)$ tend uniformly to
$y$ on that compact set as $x\to+\infty$. For all sufficiently large
$x$ they lie below $g$ there and below $A$ elsewhere. Consequently
$C=A$ on a right end interval. The residual is again zero there.
These alternatives include $F=0$, whose assertion is immediate.

On the remaining finite annulus, $C-A$ is $C^{1,1}$. More explicitly,
$|C-A|\leq\|F\|_\infty$, $|(C-A)'|\leq1$, and
$\|(C-A)''\|_{L^\infty}\leq L_2+\|A''\|_\infty$.
If $x=\log r$ is expressed in any smooth real coordinate chart on that
annulus, then almost everywhere
\[
 |\nabla^2(C-A)(x)|\leq
 (L_2+\|A''\|_\infty)|\nabla x|^2+|\nabla^2x|.
\]
Together with the exact smooth end formulas this gives a global
$C^{1,1}$ bound using $\alpha,F$ and the two end cutoffs just constructed.
No uniform bound on the location of those cutoffs is asserted for
families approaching a change in the infinity alternative.
\end{proof}

\section{Actual measures and stability}
\begin{corollary}\label{cor:mass}
The actual non-pluripolar measure of $E_F$ has mass one. For finite
$a<b$,
\[
 \MA_\kappa(E_F)(\{a<\log r<b\})=C'(b)^k-C'(a)^k.
\]
It gives zero measure to the pole and has zero density on every free
interval. For $F,G\in C_c^2(\mathbf R)$, the canonically extended
residuals satisfy
\[
 \|(E_F-\psi)-(E_G-\psi)\|_{C^0(\mathbf P^k)}
       \leq\|F-G\|_\infty.
\]
\end{corollary}
\begin{proof}
On an annulus the bounded local product is computed from the convex
total profile. For a smooth profile its Euclidean density is
$k!\pi^{-k}(C')^{k-1}C''r^{-k}$; polar integration gives
$k(C')^{k-1}C''\,dx$. Smooth convex approximation and continuity of
the bounded Bedford--Taylor product give the displayed formula.
Continuity of $C'$ excludes boundary atoms. At the pole use the strict
superlevel restrictions of $\max(E_F,-j)$; their increasing union is
the finite locus, and the non-pluripolar product gives no pole mass.
There is no hyperplane mass since the envelope is smooth at infinity.
The asymptotic slopes $0,1$ therefore give total mass one. A free
interval is affine, so its radial eigenvalue and top density are zero.
When its slope is positive the smooth Hessian has rank $k-1$.

For $\delta=\|F-G\|_\infty$, subtracting $\delta$ from any competitor
for $F$ gives a competitor for $G$, and conversely. Taking envelopes
gives the uniform bound on the finite locus. The smooth residual
extensions give the same bound at the pole and infinity. This asserts
$C^0$ stability, not convergence in $C^{1,1}$.
\end{proof}

Choose a smooth nonnegative compactly supported bump $\chi$ with
$\chi(x_0)=1$, and take $F=-M\chi$, where $M>A(x_0)-x_0$.
Then $L<0$, so \eqref{eq:tail} holds on an open right tail.
This gives an actual compact example with a retained sublogarithmic
pole, a globally $C^{1,1}$ residual, and a nonempty zero-density region.
Its regularity follows from the convex obstacle and end formulas,
not from a prescribed determinant vanishing order.

\section{Uniform families and derivative stability}
Fix a finite interval $[a,b]$, $M>0$, and the parameters $k,\alpha$.
Write $\mathcal F_{a,b,M}$ for the real $C^2$ obstacles supported in
$[a,b]$ with $\|F\|_\infty+\|F''\|_\infty\le M$, and let
$R_F$ denote the canonically extended residual $E_F-\psi$.
Norms below are taken in one fixed finite smooth projective atlas.

\begin{theorem}\label{th:uniform}
There is $K=K(k,\alpha,a,b,M)$ such that
\[
 \|R_F\|_{C^{1,1}}\le K\qquad(F\in\mathcal F_{a,b,M}).
\]
For every $0<\gamma<1$ there is $K_\gamma$, with the same fixed
parameter dependence, such that for $F,G\in\mathcal F_{a,b,M}$ and
$\delta=\|F-G\|_\infty\le1$,
\[
 \|R_F-R_G\|_{C^{1,\gamma}}
       \le K_\gamma\delta^{(1-\gamma)/2}.
\]
These bounds include changes between the infinity alternatives in
Theorem~\ref{th:main}. They assume no uniform right contact cutoff.
\end{theorem}
\begin{proof}
The left tangents of $A$ tend uniformly to $-\infty$ on $[a,b]$.
Choose $a_0<a$, depending only on $\alpha,a,b,M$, so that every
tangent at $x<a_0$ lies below $A-M$ on $[a,b]$. It is below
$A+F$ globally and touches $A$ at $x$, where $F=0$.
Thus $C_F=A$ on the same left interval for every $F$ in the family.
All residuals vanish on the same neighborhood of the pole.
The bounds
\[
 |C_F-A|\le M,\quad |C_F'-A'|\le1,\quad
 0\le C_F''\le\|A''\|_\infty+M\quad\hbox{a.e.}
\]
give uniform real Hessian bounds on any fixed compact annulus.

Put $B=\max(b,0)+1$. For $x>B$, the obstacle is $A$.
At contact points $C_F'=A'$. At a free point, the slope is
$A'(d)\ge A'(x)$ if the free interval has a finite right contact
endpoint $d$. An unbounded right free interval has slope one.
Consequently, at every $x>B$,
\[
 0\le1-C_F'(x)\le1-A'(x)\le e^{-x}.
\]
The locally Lipschitz functions $C_F',A'$ agree on the contact set.
A locally Lipschitz function has derivative zero almost everywhere
on its zero set: use its differentiability at density points and
zero-set points approaching them. Hence $C_F''=A''$ almost
everywhere on that contact set, whereas $C_F''=0$ on the free set.
The reference derivative formula now gives
\[
 0\le C_F''(x)\le A''(x)
       \le(1+\alpha(1-\alpha))e^{-x}\quad\hbox{a.e. on }x>B.
\]

In a projective chart $Z_j\ne0$, $j\ge1$, set
$\zeta=Z_0/Z_j$, $y=(Z_i/Z_j)_{i\ge1,i\ne j}$ and
\[
 H(y)=1+\sum|y_i|^2,\qquad R=|\zeta|^2/H(y),\qquad x=-\log R.
\]
On fixed relatively compact chart pieces,
$|\nabla R|\le c\sqrt R$ and $|\nabla^2R|\le c$.
Put $d_F=1-C_F'$ and express the envelope as
$\mathcal E(R)=C_F(-\log R)+\log R-\log(1+R)$ for $R>0$.
Then
\[
 \mathcal E_R=\frac{d_F}{R}-\frac1{1+R},\qquad
 \mathcal E_{RR}=\frac{C_F''-d_F}{R^2}+\frac1{(1+R)^2}
 \quad\hbox{a.e.}
\]
The preceding inequalities give $|\mathcal E_R|\le2$ and
$|\mathcal E_{RR}|\le c_\alpha/R+1$. The real chain rule therefore
bounds $\nabla^2\mathcal E$ uniformly in these projective charts.
Also $\nabla\mathcal E\to0$ as $\zeta\to0$.
The limit $L_F=\lim_{x\to\infty}(C_F(x)-x)$ lies in $[-M,0]$, and
integrating $0\le1-C_F'\le e^{-x}$ yields
$|C_F(x)-x-L_F|\le e^{-x}$.
The value error is thus $O(R)$, so the extended envelope is
differentiable at $\zeta=0$ with derivative zero.

The continuous gradient has a local $W^{1,\infty}$ bound on
coordinate boxes disjoint from $\zeta=0$, hence an actual Lipschitz
representative there by convolution. Subdivide a straight coordinate
segment into such boxes and pass to its possible single crossing of
$\zeta=0$ by gradient continuity. A segment contained in that
hyperplane has zero gradient throughout. This proves a uniform
$C^{1,1}$ bound across infinity. Subtract the fixed smooth $\psi$
there, and combine this bound with the common left cutoff and the
fixed intervening annulus to obtain the global assertion.

Finally $Z=R_F-R_G$ satisfies $\|Z\|_\infty\le\delta$ by
Corollary~\ref{cor:mass} and has a uniform Lipschitz gradient.
On a smaller fixed chart, a real directional increment $h$ gives
\[
 |\partial_e Z|\le2\delta/h+K h.
\]
Choosing $h=\sqrt{\delta/K}$ when within the chart radius, and
enlarging the fixed constant otherwise, gives
$\|\nabla Z\|_\infty\le c\sqrt\delta$.
For points separated by $t$,
\[
 |\nabla Z(x)-\nabla Z(y)|\le
             \min(c\sqrt\delta,2Kt).
\]
Divide by $t^\gamma$ and split at $t=\sqrt\delta$.
A finite atlas gives the stated $C^{1,\gamma}$ estimate.
The case $\delta=0$ follows directly from the uniform value estimate.
No $C^{1,1}$ convergence is inferred.
\end{proof}

\section{All competitors on a nonclosed product}
Let $(Y,h_0)$ be any connected compact positive Hermitian $m$-fold,
and put $X=\mathbf P^k\times Y$, $H=\kappa+h_0$,
$N=k+m$, and $V_0=\int_Yh_0^m$. Pull $\psi$ and $F$ back to $X$.
Define $E_F^H$ by \eqref{eq:envelope} with the full family of genuine
$H$-psh potentials, the pulled-back obstacle, and ordinary order.

\begin{theorem}\label{th:product}
The full product envelope is the pullback of $E_F$ and
\[
 \MA_H^{\rm ext}(E_F^H)=
       \binom Nk\,\MA_\kappa(E_F)\otimes h_0^m,
 \qquad \int_X\MA_H^{\rm ext}(E_F^H)=\binom Nk V_0.
\]
Its residual is globally $C^{1,1}$ and it belongs to the exact
stretching class $\Eenv(X,H,\psi)$. On any one fixed such product,
the uniform bounds and derivative stability of
Theorem~\ref{th:uniform} hold for the pulled-back residuals.
\end{theorem}
\begin{proof}
Every genuine $H$-psh competitor restricts on a projective slice to a
genuine $\kappa$-psh function or to $-\infty$ identically. The exact
lower Hessian in this assertion can be verified locally even when
$H$ is nonclosed. Choose a smooth quadratic correction $\rho$ with
$\D\rho\geq H$ on a coordinate polydisk. For a standard nonnegative
radial convolution kernel set
$v_\epsilon=(v+\rho)*\eta_\epsilon-\rho$ on a smaller polydisk.
The psh convolutions decrease to the canonical $v+\rho$ and
\[
 \D v_\epsilon\geq(\D\rho-H)*\eta_\epsilon-\D\rho
       \geq-H-o(1)\omega_{\rm Eucl},
\]
uniformly in the background error. Restrict this smooth inequality
to each slice. On each fixed slice the restricted psh convolutions
decrease either to $-\infty$ identically or to a psh function locally
in $L^1$; in the latter case the error tends to zero distributionally.
Subtracting $\rho$ preserves the canonical values and gives the exact
$\kappa$ lower bound. These local restrictions patch. A global slice
is consequently genuine or identically $-\infty$.

The everywhere obstacle constraint makes every nontrivial slice a
competitor for $E_F$. Thus $v\leq\operatorname{pr}_1^*E_F$ pointwise,
and upper regularization preserves this inequality since the right
side is upper semicontinuous. Conversely every base competitor lifts
because $h_0\geq0$, giving equality. No transitive group on $Y$ and
no restriction to separated competitors is used.

For bounded finite-level truncations of the pulled-back potential,
multilinearity of the locally corrected products gives only the term
with $k$ projective and $m$ factor entries. Restrict to strict
superlevel sets and let the level tend to infinity; this is the
actual non-pluripolar construction, so the same binomial identity
holds. It gives the finite mass stated. The residual bounds and
extensions pull back as well. If $K_0=\|F\|_\infty$, then
$E_F^H\geq\psi-K_0$. The genuine potential $a_1=E_F^H$ supplies
the defining inequality at $t=1$. For $t>1$, use $a_t=\psi-tK_0$, which satisfies
$t^{-1}a_t+(1-t^{-1})\psi=\psi-K_0\leq E_F^H$.
This proves the exact stretching definition, together with ordinary
order. The reference is genuine; no model-envelope identification
or unbounded-product continuity theorem has been substituted.
The uniform family and derivative estimates pull back to the fixed
smooth product atlas, with constants depending also on that atlas.
\end{proof}

For a concrete nonclosed example let
$Y=(\mathbf C^2\setminus\{0\})/\langle\zeta\mapsto e\zeta\rangle$.
With $s=\log|\zeta|^2$ and $d^c=i(\bar\partial-\partial)/(4\pi)$, put
$h_0=\D s+ds\wedge d^cs$. Both terms are scaling-invariant and descend.
The first is positive on complex tangential directions and the
second on the radial direction, so $h_0>0$.
Moreover $dh_0=-ds\wedge\D s\ne0$.
The factor is diffeomorphic to $S^3\times S^1$, with first Betti
number one, so is not K\"ahler. The product with $\mathbf P^2$ is a
compact non-K\"ahler fourfold, and $dH\ne0$. No zero-lower-mass
property of this factor is needed or asserted here.

\section{Literature and the answered scope}
Ahn's \cite[Theorem 1.4(ii)]{AhnC} proves continuity of weighted
extremals for finite continuous weights on compact Hermitian manifolds
when the unweighted extremal is continuous. His
\cite[Theorem 1.1(ii)]{AhnH} gives H\"older regularity with a uniform
density in capacity. Those theorems use weights valued in $\mathbf R$;
the weight here has a genuine $-\infty$ pole. The obstacle set here
is the entire manifold and has full capacity density. Thus this
construction does not exhibit a failure of that density assumption.
Its conclusions concern a singular unbounded weight, retained
ordinary type, and regularity of the pole-subtracted residual. The
envelope itself retains its pole and is not a globally H\"older
function. No literature-priority assertion is made.

The foundational psh envelopes and removable extensions used above
are the local statements in \cite[Chapter I, Sections 4--5]{Dem}.
All product calculations use bounded local products as in
\cite[Section 2]{BT}; smooth-background corrections and positivity
at bounded levels can be justified as in
\cite[Section 2.1, Proposition 2.1]{BGL}. The nonclosed unbounded
calculation is obtained from finite-level strict-superlevel
restrictions, not from a global theorem requiring BMP.

This gives a specific radial singular-obstacle regularity and
stability criterion with actual free-region density zeros, its
canonical extension across a nonlogarithmic pole, and an all-competitor
nonclosed product identity. Arbitrary prescribed types, general
compact obstacle sets and ideal weighted envelopes remain outside
the theorem. The ordinary model of the selecting potential is not
assumed to preserve its sublogarithmic pole.

\begin{thebibliography}{9}
\bibitem{AhnC} H. Ahn, \emph{Characterization of continuity of
Siciak--Zaharjuta extremal functions on compact Hermitian manifolds},
arXiv:2604.01885v3 (2026).
\bibitem{AhnH} H. Ahn, \emph{H\"older regularity of Siciak--Zaharjuta
extremal functions on compact Hermitian manifolds},
arXiv:2604.01887v3 (2026).
\bibitem{Dem} J.-P. Demailly, \emph{Complex Analytic and Differential
Geometry}, Chapter I.
\bibitem{BT} E. Bedford and B. A. Taylor, \emph{A new capacity for
plurisubharmonic functions}, Acta Math. 149 (1982), 1--40.
\bibitem{BGL} S. Boucksom, V. Guedj and C. H. Lu,
\emph{Volumes of Bott--Chern classes}, arXiv:2406.01090.
\end{thebibliography}
\end{document}
